\documentclass[11pt]{article}

\usepackage[margin=1in]{geometry}
\usepackage{amsmath,amssymb,amsthm,mathtools}
\usepackage{hyperref}
\usepackage{enumitem}

\hypersetup{colorlinks=true,linkcolor=blue,urlcolor=blue,citecolor=blue}
\emergencystretch=2em

\newtheorem{theorem}{Theorem}
\newtheorem{lemma}[theorem]{Lemma}
\newtheorem{corollary}[theorem]{Corollary}
\newtheorem{remark}[theorem]{Remark}

\newcommand{\T}{\mathbb T}
\newcommand{\R}{\mathbb R}
\newcommand{\N}{\mathbb N}
\newcommand{\E}{\mathbb E}
\newcommand{\dist}{\operatorname{dist}}
\newcommand{\Lip}{\operatorname{Lip}}

\title{Fixed Smooth Targets Are Poorly Approximated by Uniformly Learnable Function Classes}
\author{Run fa\_banach\_001, lane 7}
\date{August 11, 2026}

\begin{document}
\maketitle

\begin{abstract}
We answer the fixed-smoothness question in arXiv:2005.10807 in the natural
high-dimensional regime.  If the unit ball of a Banach function space $X$
admits a uniform Monte-Carlo integration rate $n^{-1/2}$, then, for $d>2k$,
the $C^k$ unit ball has $L^2$ approximation width relative to $tB_X$ bounded
below by $t^{-2k/(d-2k)}$.  A Baire-category argument produces one fixed
$C^k$ target with the corresponding limsup obstruction for squared error.
The proof uses disjoint smooth bumps which avoid the small averaging balls
needed to make the integration operators continuous on $L^2$.
\end{abstract}

\section{Statement}

Work on the flat torus $Q=\T^d$, identified with $[0,1]^d$ with periodic
boundary conditions, and let $P$ be normalized Lebesgue measure.  For an
integer $k\ge1$, equip $C^k(Q)$ with
\[
  \|g\|_{C^k}=\max_{|\alpha|\le k}\|D^\alpha g\|_\infty.
\]
Let $X$ be a Banach space continuously embedded in $L^2(Q)$.  Assume its unit
ball satisfies the uniform Monte-Carlo estimate
\begin{equation}\label{eq:MC}
 \E\sup_{\|f\|_X\le1}
 \left|\frac1n\sum_{i=1}^n f(U_i)-Pf\right|
 \le C_Xn^{-1/2},
\end{equation}
where the $U_i$ are independent and uniform on $Q$.  This is the hypothesis
used in arXiv:2005.10807; in particular it holds for its Barron and tree-like
neural-network spaces.

\begin{theorem}\label{thm:main}
Assume \eqref{eq:MC} and $d>2k$.  There are constants $c,t_0>0$ such that
\begin{equation}\label{eq:width}
 \sup_{\|g\|_{C^k}\le1}\inf_{\|f\|_X\le t}\|g-f\|_{L^2(Q)}
 \ge c\,t^{-\frac{2k}{d-2k}}
 \qquad(t\ge t_0).
\end{equation}
There also exists a single $g\in C^k(Q)$, with $\|g\|_{C^k}\le1$, such that
for every $\eta>4k/(d-2k)$,
\begin{equation}\label{eq:fixed}
 \limsup_{t\to\infty}t^\eta
 \inf_{\|f\|_X\le t}\|g-f\|_{L^2(Q)}^2=\infty.
\end{equation}
\end{theorem}

For each fixed $k$, the hypothesis $d>2k$ holds in all sufficiently large
dimensions.  Thus Theorem~\ref{thm:main} proves the source paper's conjectured
$C^k$ replacement in precisely its intended fixed-smoothness,
high-dimensional regime.

\section{Smooth empirical discrepancy}

The source asks about the integral probability metric generated by functions
which are $1$-Lipschitz and have bounded $k$-th derivative.  Modulo constants,
this class is equivalent, up to constants depending on $d,k$ and the fixed
derivative bound, to a bounded subset of $C^k(Q)$.  One direction is immediate;
the other follows from the standard finite-difference interpolation
inequalities controlling intermediate derivatives by the first and $k$-th
derivatives.  Constants do not affect integration against a signed measure of
total mass zero.

Kloeckner proved that for every probability measure $\mu$ on a bounded
$d$-dimensional domain,
\begin{equation}\label{eq:known-upper}
 \E\|\mu_n-\mu\|_{(C^k)^*}\lesssim_{d,k}
 \begin{cases}
 n^{-k/d},&d>2k,\\
 (\log n)n^{-1/2},&d=2k,\\
 n^{-1/2},&d<2k.
 \end{cases}
\end{equation}
For uniform measure in the high-dimensional regime the first rate is sharp.
We record an elementary lower bound in a form that will also survive
mollification.

For $\varepsilon>0$ and a point set $S=(x_1,\dots,x_n)\subset Q$, define
\begin{equation}\label{eq:average}
 A_{S,\varepsilon}f
 =\frac1n\sum_{i=1}^n\frac1{|B_\varepsilon|}
   \int_{B_\varepsilon(x_i)}f(y)\,dy,
 \qquad Af=Pf.
\end{equation}
The balls are periodic Euclidean balls.

\begin{lemma}[mollified smooth discrepancy]\label{lem:bump}
For every $d,k\ge1$ there are $a,c>0$ such that, with
$\varepsilon_n=a n^{-1/d}$, every $n$-point set $S\subset Q$ satisfies
\begin{equation}\label{eq:ylower}
 \|A_{S,\varepsilon_n}-A\|_{(C^k)^*}\ge c n^{-k/d}.
\end{equation}
The constants depend only on $d$ and $k$.
\end{lemma}

\begin{proof}
Fix a nonnegative $\psi\in C_c^\infty((-1/4,1/4)^d)$ with
$b:=\int\psi>0$.  Let $m=\lceil10n^{1/d}\rceil$, $h=m^{-1}$, and use the
periodic grid of $m^d$ centers $z_j$.  Take $a=1/100$.

Call $z_j$ bad if
\[
 \|z_j-x_i\|_\infty\le \varepsilon_n+h/4
\]
for some sample point $x_i$.  For a fixed $x_i$, the number of grid centers
in this cube is at most
\[
 (2m\varepsilon_n+3)^d\le4^d,
\]
because $m\varepsilon_n\le 11/100$.  Hence there are at most $4^dn$ bad
centers, while $m^d\ge10^dn$.  A fixed positive fraction of all centers is
therefore good.

Let $M_k=\max_{|\alpha|\le k}\|D^\alpha\psi\|_\infty$ and set
\[
 g_S(x)=M_k^{-1}h^k\sum_{z_j\ \mathrm{good}}
 \psi\!\left(\frac{x-z_j}{h}\right).
\]
The supports are disjoint.  For $|\alpha|\le k$,
\[
 \|D^\alpha g_S\|_\infty
 \le M_k^{-1}h^{k-|\alpha|}\|D^\alpha\psi\|_\infty\le1,
\]
so $\|g_S\|_{C^k}\le1$.  By the definition of good centers, the support of
$g_S$ misses every ball $B_{\varepsilon_n}(x_i)$, and consequently
$A_{S,\varepsilon_n}g_S=0$.  On the other hand,
\[
 Ag_S=M_k^{-1}b\,\#\{j:z_j\text{ good}\}\,h^{d+k}
 \ge c_{d,k}h^k\ge c'_{d,k}n^{-k/d}.
\]
Changing the sign of $g_S$ if necessary proves \eqref{eq:ylower}.
\end{proof}

Taking $\varepsilon=0$ in the same construction gives, for uniform $P$ and
every empirical $P_n$, the deterministic lower bound
$\|P_n-P\|_{(C^k)^*}\gtrsim n^{-k/d}$.  Together with
\eqref{eq:known-upper}, this answers the source's smooth empirical-rate
question sharply when $d>2k$.

\section{Simultaneous operator estimates}

We next choose point sets for which the lower bound of
Lemma~\ref{lem:bump} coexists with the fast estimate on $X$ and a uniform
bound on $L^2$.

\begin{lemma}\label{lem:operators}
There are point sets $S_n\subset Q$, $|S_n|=n$, for which, writing
$T_n=A_{S_n,\varepsilon_n}-A$,
\begin{align}
 \|T_n\|_{X^*}&\le C_1 n^{-1/2},\label{eq:xest}\\
 \|T_n\|_{(C^k)^*}&\ge c_1 n^{-k/d},\label{eq:yest}\\
 \|T_n\|_{(L^2)^*}&\le C_2.\label{eq:zest}
\end{align}
\end{lemma}

\begin{proof}
Choose $S=(U_1,\dots,U_n)$ randomly.  Averaging first over the ball variable
and then using translation invariance of $P$, \eqref{eq:MC} gives
\[
 \E\|A_{S,\varepsilon_n}-A\|_{X^*}\le C_Xn^{-1/2}.
\]
Notice that this does not require the $X$-norm to be translation invariant:
for every fixed shift, $(U_1+y,\dots,U_n+y)$ has the original joint law.

Let $K_\varepsilon=|B_\varepsilon|^{-1}1_{B_\varepsilon}$.  The functional
$T_n$ has the $L^2$ density
\[
 q_S(x)=\frac1n\sum_{i=1}^nK_{\varepsilon_n}(x-U_i)-1.
\]
Independence and centering yield
\[
 \E\|q_S\|_2^2
 =\frac1n\big(\|K_{\varepsilon_n}\|_2^2-1\big)
 \le \frac{1}{n|B_{\varepsilon_n}|}=C_{d,a},
\]
and hence $\E\|T_n\|_{(L^2)^*}\le C_{d,a}^{1/2}$.  Markov's inequality
shows that with positive probability both desired upper bounds hold with
three times their expectations.  The lower bound \eqref{eq:yest} holds for
every realization by Lemma~\ref{lem:bump}; selecting one realization proves
the claim.
\end{proof}

\section{Width separation and a fixed target}

We include the short separation argument to make the exponent transparent.

\begin{lemma}[functional separation]\label{lem:separation}
Let $X,Y\hookrightarrow Z$ be Banach spaces.  Suppose $T_n\in Z^*$ satisfy
\[
 \|T_n\|_{X^*}\le C_Xn^{-\alpha},\qquad
 \|T_n\|_{Y^*}\ge c_Yn^{-\beta},\qquad
 \|T_n\|_{Z^*}\le C_Z
\]
with $0<\beta<\alpha$.  Then, for all large $t$,
\[
 \rho(t):=\sup_{\|y\|_Y\le1}\dist_Z(y,tB_X)
 \ge c\,t^{-\beta/(\alpha-\beta)}.
\]
\end{lemma}

\begin{proof}
Choose $y_n\in B_Y$ with $|T_ny_n|\ge(c_Y/2)n^{-\beta}$.  For
$x\in tB_X$,
\[
 (c_Y/2)n^{-\beta}
 \le C_Z\|y_n-x\|_Z+C_Xt n^{-\alpha}.
\]
Choose an integer $n\asymp t^{1/(\alpha-\beta)}$ large enough that the last
term is at most $(c_Y/4)n^{-\beta}$.  Taking the infimum over $x$ gives the
claim.
\end{proof}

The source's constructive fixed-point lemma imposes the stronger condition
$\beta<\alpha/2$.  For existence of a fixed hard target, Baire category
removes this loss.

\begin{lemma}[Baire fixed-target upgrade]\label{lem:baire}
Let $X,Y\hookrightarrow Z$ be Banach spaces, and suppose
\[
 \rho(t)=\sup_{\|y\|_Y\le1}\dist_Z(y,tB_X)\ge c t^{-q}
\]
for all large $t$.  Then for every $\gamma>q$ there exists $y\in B_Y$ such
that
\[
 \limsup_{t\to\infty}t^\gamma\dist_Z(y,tB_X)=\infty.
\]
In fact, one $y$ may be chosen to have this property simultaneously for all
$\gamma>q$.
\end{lemma}

\begin{proof}
Fix $\gamma>q$ and write $E_t(y)=\dist_Z(y,tB_X)$.  Suppose every $y\in Y$
satisfies $\sup_{t\ge1}t^\gamma E_t(y)<\infty$.  Then
\[
 Y=\bigcup_{M=1}^\infty F_M,
 \qquad
 F_M=\{y:E_t(y)\le Mt^{-\gamma}\text{ for all }t\ge1\}.
\]
Each $F_M$ is closed because $E_t$ is continuous on $Y$.  By Baire's theorem,
some $F_M$ contains a ball $y_0+rB_Y$.  Put $\delta=r/2$.  For every
$u\in B_Y$, both $y_0$ and $y_0+\delta u$ lie in $F_M$.  Subtracting
$tB_X$-approximants to these two points gives
\[
 E_{2t}(\delta u)\le2Mt^{-\gamma}.
\]
The exact scaling identity
$E_{2t}(\delta u)=\delta E_{2t/\delta}(u)$ therefore implies
\[
 \rho(s)\le C_{M,r,\gamma}s^{-\gamma}
\]
for all large $s$, contradicting $\rho(s)\ge cs^{-q}$.

Thus for each rational $\gamma>q$, the set of points with unbounded
$t^\gamma E_t(y)$ is dense $G_\delta$ by the same argument applied in every
open ball.  Intersecting over rational $\gamma>q$ and normalizing a nonzero
point gives a single $y\in B_Y$ working for every real $\gamma>q$.
\end{proof}

\begin{proof}[Proof of Theorem~\ref{thm:main}]
Apply Lemma~\ref{lem:separation} to the operators of
Lemma~\ref{lem:operators} with
\[
 Y=C^k(Q),\qquad Z=L^2(Q),\qquad
 \alpha=\frac12,\qquad\beta=\frac{k}{d}.
\]
The condition $d>2k$ is exactly $\beta<\alpha$, and
\[
 \frac{\beta}{\alpha-\beta}
 =\frac{k/d}{1/2-k/d}=\frac{2k}{d-2k}.
\]
This proves \eqref{eq:width}.  Lemma~\ref{lem:baire} gives one
$g\in B_{C^k}$ for which
\[
 \limsup_{t\to\infty}t^\gamma
 \inf_{\|f\|_X\le t}\|g-f\|_2=\infty
 \quad\text{for every }\gamma>\frac{2k}{d-2k}.
\]
Squaring and putting $\eta=2\gamma$ proves \eqref{eq:fixed}.
\end{proof}

\section{Boundary of the result}

The exponent transition has a structural meaning.  When $d>2k$, smooth
empirical integration is of order $n^{-k/d}$, strictly slower than the
$n^{-1/2}$ rate on $X$, and the separation principle applies.  When $d<2k$,
the $C^k$ empirical rate is already $n^{-1/2}$; at $d=2k$ it is parametric up
to a critical logarithm.  Thus the present mechanism has no exponent gap at
$d\le2k$.  This is consistent with the source's formulation: $k$ is fixed
while dimension becomes large.

The theorem is abstract in $X$.  Substituting the Barron space or the
tree-like spaces from arXiv:2005.10807 gives the claimed neural-network
consequence immediately.  The result does not classify which smooth targets
are hard, and it does not address loss functionals such as logistic loss.

\section{References}

\begin{itemize}[leftmargin=*]
\item W. E and S. Wojtowytsch, \emph{Kolmogorov Width Decay and Poor
Approximators in Machine Learning: Shallow Neural Networks, Random Feature
Models and Neural Tangent Kernels}, Research in the Mathematical Sciences 8,
5 (2021); \href{https://arxiv.org/abs/2005.10807}{arXiv:2005.10807}.
\item B. Kloeckner, \emph{Empirical measures: regularity is a counter-curse
to dimensionality}, ESAIM: Probability and Statistics 24 (2020), 408--434;
\href{https://arxiv.org/abs/1802.04038}{arXiv:1802.04038}.
\end{itemize}

\end{document}
